\documentclass[11pt]{letter}
\usepackage[margin=1in]{geometry}

\signature{Decory Edwards}

\begin{document}

\begin{letter}{Hiring Committee \\
Federal Reserve Bank of Atlanta}

\opening{Dear Hiring Committee,}

I am writing to apply for the Summer 2026 Research Internship at the Federal Reserve Bank of Atlanta. I am a Ph.D. candidate in Economics at Johns Hopkins University, expecting to complete my degree in May 2026. My research focuses on macroeconomics and household finance, and I would greatly value the opportunity to advance my work while engaging with the Atlanta Fed’s research community.

My job market paper develops a heterogeneous-agent life-cycle model to study wealth accumulation and inequality, incorporating heterogeneity in returns and income risk. The project combines quantitative modeling with empirical discipline and is designed to generate insights relevant to macroeconomic policy and distributional outcomes. In parallel work, I use micro and panel data from the Health and Retirement Study to study the relationship between trust and financial outcomes, applying econometric methods suited for panel data and causal identification. Together, these projects reflect my broader interest in understanding how micro-level behavior aggregates into macroeconomic outcomes.

During the internship, there are a number of directions I could take my current research agenda. The current of the model in my job market paper does not include bequest motives, portfolio choice, or overlapping generations. Each of these modeling choices, thoguh they would add to computational complexity, are feautures of models in the heterogeneous agent framework which matter for the resulting model distribution of wealth.

I would also be interested in focusing on the project establishing an empirical relationship between trust and returns in the HRS. There are many directions this paper could be taken. An obvious one is that I've been using the RAND longitudinal version of the HRS data. Although this has been useful, it is the best of a worst case situation: computing returns naturally requires sensitive information about respondents, so I should likely be using some data which requires more clearance to use. The most important thing in this direction would be finding a dataset that has both a measure of trust and the senstitive information on changes in individuals' asset holdings required to compute returns.

I am very interested in the Atlanta Fed’s collaborative research environment and its focus on policy-relevant economic analysis. I would be excited to contribute to and learn from the Department’s seminar series and ongoing research initiatives.

Thank you for your time and consideration.

\closing{Sincerely,}

\end{letter}

\end{document}